To systematically evaluate the performance impact of garbage collection across different workloads and platforms, a modular and reproducible experimental framework was designed and implemented. This framework provides the necessary infrastructure for compiling source code, running benchmarks consistently, and collecting performance data.

\subsection{Framework Architecture}

The benchmark system is organized within a top-level directory named `Benchmark/`. This directory contains the core automation scripts, specific benchmark implementations categorized by workload type, and tooling for data analysis and visualization.

\subsubsection{Directory Structure}
The framework employs a hierarchical structure to organize benchmark tasks:
\begin{itemize}
    \item \textbf{Benchmark/}: The root directory containing automation scripts and subdirectories for benchmark categories.
    \item \textbf{Compute/}: Contains subdirectories for CPU-intensive tasks. Each subdirectory (e.g., `FibonacciLoop/`, `Matrix/`, `Pi/`) holds the language-specific implementations (Go, Java, Rust) for that compute task.
    \item \textbf{IO/}: Contains subdirectories for I/O-bound tasks, such as `LargeFile/` for file operations and `WebServer/` for simulating network request handling. Similar to `Compute/`, each task directory contains the respective language implementations.
\end{itemize}
This modular organization facilitates the addition of new benchmark tasks or language implementations. A visual representation of the structure is shown in Figure \ref{fig:framework-structure}.

\begin{figure}[h!]
\centering
\begin{tikzpicture}[
  node distance=1cm and 1.2cm,
  folder/.style={draw, thick, minimum width=2.6cm, minimum height=1cm, text centered, rounded corners},
  script/.style={draw, minimum width=2.6cm, minimum height=0.8cm, text centered},
  line/.style={draw, -{Latex[length=2mm]}}
]

% Top node
\node (benchmark) [folder] {Benchmark/};

% First layer scripts
\node (setup)   [script, below left=of benchmark] {setup.sh};
\node (prepare) [script, below right=of benchmark] {prepare.sh};

% Second layer: benchmark & clean scripts
\node (benchsh) [script, below=of setup] {benchmark.sh};
\node (cleansh) [script, below=of prepare] {clean.sh};

% graph.py
\node (graph) [script, below=of benchsh] {graph.py};

% Compute & IO folders
\node (compute) [folder, below left=2.4cm and -0.2cm of graph] {Compute/};
\node (io)      [folder, below right=2.4cm and -0.2cm of graph] {IO/};

% Subfolders in Compute and IO
\node (fib) [folder, below=of compute] {FibonacciLoop/};
\node (web) [folder, below=of io] {WebServer/};

\node (fibdots) [below=0.2cm of fib] {...};
\node (webdots) [below=0.2cm of web] {...};

% Connections
\draw [line] (benchmark) -- (setup);
\draw [line] (benchmark) -- (prepare);
\draw [line] (benchmark) -- (benchsh);
\draw [line] (benchmark) -- (cleansh);
\draw [line] (benchmark) -- (graph);
\draw [line] (benchmark) -- (compute);
\draw [line] (benchmark) -- (io);
\draw [line] (compute) -- (fib);
\draw [line] (io) -- (web);

\end{tikzpicture}
\caption{High-level structure of the benchmark framework.}
\label{fig:framework-structure}
\end{figure}

\subsubsection{Automation Scripts}
A set of shell scripts automates the benchmarking workflow, ensuring consistency and reproducibility:
\begin{itemize}
    \item \texttt{setup.sh}: Prepares the testing environment by installing necessary dependencies such as the Rust compiler, Go compiler, Java Development Kit (JDK), and the `hyperfine` benchmarking tool. It includes OS detection for compatibility with macOS and Linux.
    \item \texttt{prepare.sh}: Compiles all Go, Java, and Rust source files found within the task subdirectories, creating executable binaries or class files. Optimization flags are used for compiled languages (e.g., `rustc -O` or `rustc -C opt-level=3`).
    \item \texttt{benchmark.sh}: Executes the core benchmarking logic for a specified task directory. It utilizes the `hyperfine` tool to run the compiled programs, collect timing data, perform statistical analysis (mean, standard deviation), handle warmup runs, and export results to CSV and Markdown formats.
    \item \texttt{run.sh}: A helper script, typically invoked by `hyperfine` within `benchmark.sh`, responsible for executing a specific language's compiled program for a given task.
    \item \texttt{clean.sh}: Removes compiled artifacts (executables, `.class` files) to ensure a clean state before subsequent runs.
    \item \texttt{graph.py}: A Python script that reads the CSV output generated by `hyperfine` and uses `matplotlib` and `pandas` to create bar charts visualizing the mean execution times for each language within a benchmark task, saving them as PNG images.
\end{itemize}

\subsection{Benchmarking Tool: Hyperfine}

The framework relies on \texttt{hyperfine}, a command-line benchmarking tool chosen for its robustness and features suitable for this research. Key capabilities leveraged include:
\begin{itemize}
    \item \textbf{Statistical Analysis}: Automatically computes mean, median, standard deviation, minimum, and maximum execution times over multiple runs.
    \item \textbf{Warmup Runs}: Allows specifying a number of initial runs (\texttt{--warmup N}) whose results are discarded, mitigating skewed measurements due to initial setup costs like JIT compilation or cache population.
    \item \textbf{Run Count}: Enables setting the number of measured runs (\texttt{--runs N}) for statistical significance.
    \item \textbf{Output Formats}: Exports results in various formats, including CSV for data analysis and Markdown for easy reporting.
    \item \textbf{Inter-process Measurement}: Accurately measures the wall-clock time of external commands.
\end{itemize}
A typical invocation within the framework looks like this:
\begin{lstlisting}[language=bash, caption={Example hyperfine usage in benchmark.sh }, style=shstyle]
hyperfine \
  --warmup 3 \
  --runs 10 \
  --export-csv "benchmark_TASK.csv" \
  --export-markdown "benchmark_TASK.md" \
  "./Rust_TASK <args>" \
  "./Go_TASK <args>" \
  "java Java_TASK <args>" # Example, actual execution might use run.sh
\end{lstlisting}

\subsection{Cross-Platform Testing} % Retained from original, slightly rephrased

The framework is designed to be executed on diverse hardware architectures to assess how performance characteristics, particularly those related to GC, might vary:
\begin{itemize}
    \item Personal computer (e.g., Macbook Pro) 
    \item Cloud-based server (e.g., Linux server, x86 architecture) 
    \item Embedded platform (e.g., ESP32, representing resource-constrained devices) 
\end{itemize}
This multi-platform testing strategy aims to provide a broader understanding of GC performance implications across different computational scales and constraints.
